\documentclass[10pt]{article}
\usepackage[latin9]{inputenc}
\usepackage[a4paper]{geometry}
\geometry{verbose,tmargin=2cm,bmargin=2cm,lmargin=2cm,rmargin=2cm}
\usepackage{xcolor}
\usepackage{float}
\usepackage{calc}
\usepackage{bm}
\usepackage{amsmath}
\usepackage{amssymb}
\usepackage{graphicx}
\PassOptionsToPackage{normalem}{ulem}
\usepackage{ulem}
\usepackage{amsthm}
\usepackage{url}  % 或 hyperref，它也包含 url 功能
\usepackage{hyperref}

\makeatletter
%%%%%%%%%%%%%%%%%%%%%%%%%%%%%% User specified LaTeX commands.
%!TEX program = xelatex
\usepackage{mathrsfs}
\usepackage{titlesec}
\usepackage{xcolor}
\usepackage{enumerate}
\usepackage{bm}
\usepackage{tikz}
\usepackage{listings}
\usetikzlibrary{arrows}
\usepackage{subfigure}
\usepackage{graphicx}
\usepackage{booktabs}\usepackage{multirow}\usepackage{upquote}
\usepackage{float}
\usepackage{pdfpages}

% additional commands
\newcommand{\mbf}[1]{\mathbf{#1}}
\newcommand{\gauss}{\mathcal{N}}
\newcommand{\set}[1]{\left\{#1\right\}}
\newcommand{\mat}[1]{\mathbf{#1}}

\lstset{language=Matlab}
\lstset{breaklines}

\input defs.tex
\newtheorem{proposition}{Proposition}\newtheorem{remark}{Remark}
\newcommand{\vectdot}[1]{\dot{\mathbf{#1}}}
\titleformat*{\section}{\centering\LARGE\scshape}
\renewcommand{\thesection}{\Roman{section}}
\lstset{language=Matlab,tabsize=4,frame=shadowbox,basicstyle=\footnotesize,
keywordstyle=\color{blue!90}\bfseries,breaklines=true,commentstyle=\color[RGB]{50,50,50},stringstyle=\ttfamily,numbers=left,numberstyle=\tiny,
  numberstyle={\color[RGB]{192,92,92}\tiny},backgroundcolor=\color[RGB]{245,245,244},inputpath=code}

\makeatother

\begin{document}
\title{CS182: Introduction to Machine Learning, Fall 2025\\
 Homework 2}
% \author{{\small (Due on Sunday, Nov. 9 at 11:59 pm (CST))}}
\author{\small{(Due on Friday, Nov. 14 at 11:59 pm (CST))}}
\date{\date{}}

\maketitle
{\small{}%
\noindent\fcolorbox{black}{white}{\begin{minipage}[t]{1\columnwidth - 2\fboxsep - 2\fboxrule}%
Notice:
\begin{itemize}
\item Please submit your assignments via Gradescope. The entry code is \uline{428YBY}.
\item Please make sure you select your answer to the corresponding question
when submitting your assignments.
\item Each person has a total of five days to be late without penalty for
all the assignments. Each late delivery less than one day will be
counted as one day.
\end{itemize}
%
\end{minipage}}}{\small\par}
\begin{enumerate}
\item \defpoints{20} {[}\textit{Parameter Estimation}{]} A 2-axis force sensor, mounted on a robot fingertip, measures an applied force $\mbf{F} = (F_x, F_y)$. Due to manufacturing tolerances and mechanical crosstalk, the sensor's voltage output $V_x$ is linearly dependent on both $F_x$ and $F_y$.

We use a Multivariate Linear Regression Model to calibrate the $V_x$ channel. This model includes a bias term $b$ and two gain terms, $m_{11}$ and $m_{12}$:
$$V_x = m_{11} F_x + m_{12} F_y + b$$

We assume the sensor's measurement noise $\epsilon$ is Gaussian, $\epsilon \sim \gauss(0, \sigma^2)$. Therefore, for the $i$-th measurement, the observation model is:
$$V_x^i = (m_{11} F_x^i + m_{12} F_y^i + b) + \epsilon_i$$

The parameter vector we want to estimate is $\mbf{w} = \begin{bmatrix} m_{11} \\ m_{12} \\ b \end{bmatrix}$.



\begin{enumerate}
    \item Given a set of $N$ i.i.d. calibration data points $\set{X} = \{ (F_x^i, F_y^i), V_x^i \}_{i=1}^N$. Write the Log-Likelihood Function $\mathcal{L}(\mbf{w} | \set{X})$ for the parameters $\mbf{w}$. \defpoints{4}

    \item We know that maximizing $\mathcal{L}(\mbf{w} | \set{X})$ is equivalent to minimizing the Sum of Squared Errors (SSE) $E(\mbf{w}) = \sum_{i=1}^{N} (V_x^i - (m_{11} F_x^i + m_{12} F_y^i + b))^2$.\defpoints{8}
    
    \begin{itemize}
        \item Write the SSE $E(\mbf{w})$ in matrix-vector form: $E(\mbf{w}) = \|\mbf{r} - \mat{X}\mbf{w}\|^2$. Clearly define the design matrix $\mat{X}$, the parameter vector $\mbf{w}$, and the observation vector $\mbf{r}$ based on the model.
        
        \item By computing the gradient of $E(\mbf{w})$ with respect to the vector $\mbf{w}$ (i.e., $\nabla_{\mbf{w}} E(\mbf{w})$) and setting it to zero, derive the Normal Equations:
        $(\mat{X}^T \mat{X}) \hat{\mbf{w}} = \mat{X}^T \mbf{r}$,
        
        % (\textbf{Hint:} You may use the matrix calculus identity $\nabla_{\mbf{x}} \|\mbf{y} - \mat{A}\mbf{x}\|^2 = -2\mat{A}^T(\mbf{y} - \mat{A}\mbf{x})$ or derive it from the expanded form $E(\mbf{w}) = (\mbf{r} - \mat{X}\mbf{w})^T(\mbf{r} - \mat{X}\mbf{w})$.)
    \end{itemize}
    
    \item During an experiment, $N=4$ known forces were applied, and the corresponding $V_x$ voltages were measured.\defpoints{8}

\begin{center}
\begin{tabular}{c c c}
\toprule
Measurement $i$ & Applied Force $(F_x, F_y)$ (N) & Measured Voltage $V_x$ (V) \\
\midrule
1 & (1, 0) & 2.1 \\
2 & (0, 1) & 0.9 \\
3 & (1, 1) & 3.0 \\
4 & (2, 0) & 4.1 \\
\bottomrule
\end{tabular}
\end{center}
\begin{itemize}
    \item Using the data provided in the table, first compute the numerical values for $\mat{X}^T \mat{X}$ and $\mat{X}^T \mbf{r}$. Then, solve the system to find the optimal parameter estimates $\hat{\mbf{w}} = (\hat{m}_{11}, \hat{m}_{12}, \hat{b})^T$. 
    \item A unique least-squares solution $\hat{\mbf{w}}$ only exists if the applied calibration forces satisfy a specific condition. From an experimental design perspective, what \textbf{physical requirement} must the set of $N=4$ applied forces $\{ (F_x^i, F_y^i) \}$ meet to ensure that the parameters $(m_{11}, m_{12}, b)$ can be uniquely determined? Briefly explain why. 
\end{itemize}
\end{enumerate}

\newpage

% \begin{enumerate}
\item \defpoints{20} {[}\textit{Parameter Estimation}{]} 
Assume that the prior distribution of $\theta$ follows a Gaussian distribution:
$$p(\theta) = \mathcal{N}(0, \sigma^2),$$
where $\theta$ is an unknown parameter and $\sigma$  is known. The conditional distribution of $x$ given $\theta$ is
$$p(x \mid \theta) = \mathcal{N}(\theta, 1).$$
Please provide detailed derivations for the following:
\begin{itemize}
	\item[(a)] Derive the marginal distribution $p(x)$. (Hint: $\int_{-\infty}^{\infty} e^{-a x^2 + b x + c} dx = \sqrt{\frac{\pi}{a}}e^{\frac{b^2}{4a} + c}$.) \defpoints{6}
    
	\item[(b)] Compute the posterior expectation $\mathbb{E}(\theta \mid x)$. \defpoints{6}
    
	\item[(c)] Suppose the prior $p(\theta)$ is unknown, and we only know that $p(x \mid \theta) = \mathcal{N}(\theta, \sigma^2).$ Find the posterior expectation $\mathbb{E}(\theta \mid x)$. \defpoints{8}
\end{itemize}


\newpage{}




\item\defpoints{20} {[}\textit{Parameter Estimation}{]}
\begin{itemize}
In class, we have learned that the MMSE is achieved when using the conditional expectation \(\E[X | \mathbf{Y}]\), i.e., \(\arg\min_{\hat{X}}\E[|X-\hat{X}(\mathbf{Y})|^2] = E[X|\mathbf{Y}]\), where \(X\) is some unknown random variable and \(\mathbf{Y}\) is the observation. However, \(\E[X|\mathbf{Y}]\) can be intractable to compute. Now, let's focus on linear estimation and minimize MSE over all linear estimators.
\begin{definition}[LLSE]
    \[\mathbb{L}[X|\mathbf{Y}] = \arg\min_{\text{linear }\hat{X}}\E[|X-\hat{X}(\mathbf{Y})|^2]\]
\end{definition}
where linear estimators \(\hat{X}(\mathbf{Y})\) are of the form 
\[\hat{X}(\mathbf{Y}) = a+\sum_{i = 1}^nb_iY_i\] for \(a,b_i\in \mathbb{R}\) where \(\mathbf{Y} = (Y_1, \cdots , Y_n)\) is a vector of observations.

\item[(a)] 
Define
\begin{align*}
    \Sigma_{\mathbf{XY}} &= \E[(X-\mu_x)(\mathbf{Y}-\mu_\mathbf{Y})^T]&\text{(cross-covariance matrix)}\\
    \Sigma_{\mathbf{Y}} &= \E[(\mathbf{Y}-\mu_\mathbf{Y})(\mathbf{Y}-\mu_\mathbf{Y})^T]&\text{(covariance matrix)}
\end{align*}
where \(\mu_X=\E(X),\mu_Y= \E[Y] = [\E[Y_1],\cdots,\E[Y_n]]\)\\
Show that 
\begin{align}
    \mathbb{L}[X|\mathbf{Y}] = \mu_X+\Sigma_{\mathbf{XY}}\Sigma_{\mathbf{Y}}^{-1}(\mathbf{Y}-\mu_\mathbf{Y})
\end{align}\defpoints{5}

\item[(b)] Show that a linear estimator \(\hat{X}(\mathbf{Y})\) achieves minimum mean square error(over all linear estimators) if and only if 
\begin{align}
    &\E[\hat{X}(\mathbf{Y})-X] = 0\\
    &\E[(\hat{X}(\mathbf{Y})-X)Y_i] = 0
\end{align}
This is called the \textbf{orthogonality principle} since it uniquely characterizes \(\mathbb{L}[X|\mathbf{Y}]\).\\
\defpoints{10}
\item[(c)] In the real world, data is observed sequentially. How do we efficiently update our LLS estimate on the arrival of new observations? For motivation, start with a simple setting. w.l.o.g., assume \(\E[X] = 0\), Suppose observations \(Y_1,Y_2,\cdots \) are orthogonal and uncorrelated, i.e., \(\langle Y_i,Y_j\rangle=\E[Y_iY_j]=\Cov(Y_i,Y_j) = 0\) for \(i\neq j\) to each other.\\
For convenience, define
\begin{align*}
    \mathbf{Y}^n&=(Y_1,\cdots,Y_n)\\
    \mu_i  &= \E[Y_i]
\end{align*}
Show that
\begin{align}
     \mathbb{L}[X|\mathbf{Y}^{n+1}]&=  \mathbb{L}[X|\mathbf{Y}^n]+ \mathbb{L}[X|Y_{n+1}]\\
     &= \mathbb{L}[X|\mathbf{Y}^n]+ \frac{\Cov(X,Y_{n+1})}{\Var(Y_{n+1})}(Y_{n+1}-\mu_{n+1})
\end{align}
\defpoints{5}\\
(Hint: Just check the orthogonality principle!)
\end{itemize}

\newpage


\newpage






\item \defpoints{20} {[}\textit{Linear Discriminant Analysis}{]} We have
a two-class dataset with label $\mathcal{Y}=\left\{ -1,+1\right\} $
and the feature vector $\bm{x}\in\mathbb{R}^{2}$. We find the means
of these two classes are $\hat{\mathbf{\bm{\mu}}}_{-1}=\left(-1,-1\right)^{\top}$
and $\hat{\mathbf{\bm{\mu}}}_{+1}=\left(1,1\right)^{\top}$, respectively.
And the estimated prior probabilities for these two class are equal
(i.e., $\hat{\pi}_{+1}=\hat{\pi}_{-1}=\frac{1}{2}$).\\
 Some additional explanations on question 4. The Fisher's linear discriminant is a supervised dimensionality
reduction method. The terms Fisher's linear discriminant and LDA are
often interchangeably. The  Fisher's linear discriminant maximizes the ratio between $\bm{S}_{B}$
and $\bm{S}_{W}$, where $\bm{S}_{B}$ denotes the between-class scatter
matrix and $\bm{S}_{W}$ denotes the within-class scatter matrix, and this
method is based on the so-called Fisher's linear discriminant rule.
However, there is another type of LDA, which is based on Bayes discriminant
rule, is presented to you in this question. By assuming $p\left(\bm{x}\left|C_{i}\right.\right)\sim\mathcal{N}\left(\bm{\mu}_{i},\bm{\Sigma}\right)$
(i.e., homoscedasticity), this type of LDA assigns $\bm{x}$ to the
group that maximizes $\pi_{i}p\left(\bm{x}\left|C_{i}\right.\right)$,
where $\pi_{i}$ represents the prior probability of each classification
$C_{i}$, and $p\left(\bm{x}\left|C_{i}\right.\right)$ represent
the $i$-th density. In question 4(a), we want you to use this type
of LDA based on Bayes discriminant rule. In question 4(b), without
homoscedasticity, the problem converts to quadratic discriminant analysis
(QDA). You can get more detailed information through:  \url{https://harvard-iacs.github.io/2018-CS109A/a-sections/a-section-6/presentation/a_section6.pdf}\\

\begin{itemize}
\item[(a)]  Assuming the covariance martix can be estimated as
\[
\hat{\mathbf{\Sigma}}=\begin{pmatrix}1 & \rho\\
\rho & 1
\end{pmatrix}
\]
with $0\leq\rho<1$ when applying LDA. Find the decision boundary
as a function of $\rho$. \defpoints{7}
\item[(b)]  Now suppose that we model each class with its own covariance matrix.
The covariance matrix of class $+1$ is
\[
\hat{\mathbf{\Sigma}}_{+1}=\begin{pmatrix}1/5 & 0\\
0 & 5
\end{pmatrix}.
\]
And the covariance matrix of class $-1$ is
\[
\hat{\mathbf{\Sigma}}_{-1}=\begin{pmatrix}5 & 0\\
0 & 1/5
\end{pmatrix}.
\]
Describe the decision rule. \defpoints{8}
\item[(c)]  When $p\left(\bm{x}\left|C_{i}\right.\right)\sim\mathcal{N}\left(\bm{\mu}_{i},\hat{\bm{\Sigma}}\right)$
(i.e., homoscedasticity), using Fisher's linear discriminant, prove that the vector $\mathbf{w}$ on which the data $\bm{x}$
are projected is $\mathbf{w}=\hat{\bm{\Sigma}}^{-1}\left(\bm{\mu}_{+1}-\bm{\mu}_{-1}\right)$.
\defpoints{5}
\end{itemize}

% \end{itemize}
\newpage

\item \defpoints{20} [\textit{Logistic Discrimination}]
Use a quadratic discriminant classifier to classify the $K$ classes $C_{1},\dots,C_{K}$. Assume that the class-conditional distributions are multivariate normal:
\[
p\left(\mathbf{x} \mid C_i\right) = \mathcal{N}\left(\mathbf{x} \mid \boldsymbol{\mu}_i, \boldsymbol{\Sigma}_i\right)
\]
where $\boldsymbol{\mu}_i$ is the mean vector and $\boldsymbol{\Sigma}_i$ is the covariance matrix for class $C_i$.

The quadratic discriminant function is given by:
\begin{equation}
	\label{eq:quad-discriminant}
	\log\frac{p\left(\mathbf{x}\mid C_{i}\right)}{p\left(\mathbf{x}\mid C_{K}\right)}=\mathbf{x}^{\top}\mathbf{W}_{i}\mathbf{x}+\mathbf{w}_{i}^{\top}\mathbf{x}+w_{i0}
\end{equation}
where $\mathbf{W}_{i} \in \mathbb{R}^{d\times d}$ is a symmetric matrix related to the covariance matrices, $\mathbf{w}_{i} \in \mathbb{R}^{d}$ is a vector, and $w_{i0} \in \mathbb{R}$ is a scalar.

\begin{itemize}
	\item[(a)] Show the relationship between the parameters $\left(\mathbf{W}_i, \mathbf{w}_i, w_{i0}\right)$ and the Gaussian parameters $\left(\boldsymbol{\mu}_i, \boldsymbol{\Sigma}_i\right)$. \defpoints{5}
	
	\item[(b)] Using the quadratic discriminants in equation \eqref{eq:quad-discriminant}, derive the expression for the posterior probability $P\left(C_i \mid \mathbf{x}\right)$ for each class $i=1,\dots,K$. Your answer should involve a softmax function. \defpoints{5}
	
	\item[(c)] Show that if all classes share the same covariance matrix, i.e., $\boldsymbol{\Sigma}_i = \boldsymbol{\Sigma}$ for all $i$, then the decision boundaries become linear. Also, find the expressions for the parameters $\mathbf{W}_{i}$, $\mathbf{w}_{i}$, and $w_{i0}$. \defpoints{5}
	
	\item[(d)] Discuss the advantages and disadvantages of quadratic discriminant compared to linear discriminant. Consider aspects such as model flexibility, number of parameters, and risk of overfitting. \defpoints{5}
\end{itemize}




\end{enumerate}
\end{document}
